\section{Introduction}
\label{sec:intro}

Recommenders are usually trained and evaluated on logs written by a production policy. Such a log records which
items the policy chose to show to which users, and whether those users clicked. A model fitted to this record
learns the response of the audiences that the policy selected, and a model that ranks well on a held-out part of
the same record is rewarded partly for reproducing the policy. The question that matters for an item is different:
how a user who was not selected for the item would respond to it. Randomised exposure measures this target
directly. In the KuaiRand logs \citep{gao2022kuairand}, a feed slot is replaced with small probability by a video
drawn uniformly from a pool of 7,583 videos, so the inserted candidates reach users whatever the policy would have
predicted.
% source: results/v11_probes/diag_gap.txt (header: items with any 7583)

We first ask where the remaining accuracy of a strong graph recommender lies on this target. The headroom is at the
item level. A leave-one-user-out item oracle, which scores each candidate by the uniform click rate of the item among
the other test users, reaches nDCG@10 of 0.5068 against 0.4735 for LEAF-LGN, a gap of $+33.31\times10^{-3}$. The
oracle reads test labels, so it is a diagnostic and not a method. The gap is concentrated in the items with the
highest uniform click rates. The post-cut log predicts the uniform rate no better than the pre-cut log (Spearman
0.405 against 0.431), so temporal drift is not the bottleneck. Personalisation fitted on uniform in-pool data adds
nothing once item means are known. What the policy log lacks is a measurement of how appealing an item is to an
audience the policy did not select.
% sources: results/v11_probes/diag_gap.txt (oracle 0.5068, gap +33.31e-3, LEAF-LGN 0.4735, Spearman rows
% post_clk +0.405 and pre_clk +0.431, quintile share in row q bin 4);
% results/v9_probes/headroom_kuairand_pure.txt (H2: item-only 0.5850, rank-64 residual 0.5851)

We address this with two signals that are estimated from the log alone and added to a graph backbone. The first is a
lifelong, exposure-conditional affinity (LEAF). For each user, LEAF uses the impressions of videos outside the test
pool before the cut-off, and summarises, for every content key of a pool video (tag, author, music), how often the
user clicked after an impression. Because the statistics are conditioned on the impression, they measure the response
of the user given exposure, not how often the policy chose the key. The second is detection-adjusted attractiveness
(DAX). DAX models each logged impression as an examination event, whose probability depends on the user and the
scenario, followed by a click that depends on the attractiveness of the item. A non-click is then weak evidence
against an item when the user rarely examines, and strong evidence when examination is likely. The attractiveness is
fitted by expectation-maximisation on the pre-cut log and shrunk towards the mean of its content tag.

We make four contributions. Each is stated with the number that supports it. Section~\ref{sec:results} gives the
full tables, and each number is traced to a result file in the replication materials.

\begin{description}
\item[C1 (headroom diagnosis).] On randomised insertions, the item oracle exceeds LEAF-LGN by
  $+33.31\times10^{-3}$ nDCG@10. This is a diagnostic that uses test labels and is never a selection input.
% source: results/v11_probes/diag_gap.txt
\item[C2 (LEAF).] Lifelong exposure-conditional channels built from out-of-pool videos transfer through content keys
  to a disjoint, randomly exposed pool. Added to LightGCN they give $+15.46$ [1.53] $\times10^{-3}$ nDCG@10, and a
  control that shuffles the users' lifelong tables removes the lifelong gain. LEAF alone does not pass the pre-declared
  Stage 1 gate against the QP family (Section~\ref{sec:results}).
% source: results/v10/stage1.txt Sections 3 and 5 (LEAF-LGN vs LightGCN; shuffle control, closed-form variant)
\item[C3 (DAX).] DAX is complementary to LEAF. On top of LEAF-LGN it gives $+11.96$ [1.47] $\times10^{-3}$ nDCG@10,
  and the same post-hoc module is applied to EASE and to the QP family. The DAX weight is selected on validation
  together with the backbone epoch and the DAX configuration; it is not learned jointly.
% source: results/v11/stageB.txt Section 2 (LEAF-LGN+DAX - LEAF-LGN); post-hoc rule from results/v11/dax_lgn_post.json
\item[C4 (main result).] Under a temporal protocol with five user-level splits, LEAF-LGN+DAX reaches nDCG@10 of
  0.4854 against 0.4634 for the QP family, a paired difference of $+22.32$ [1.83] $\times10^{-3}$ with a positive
  difference in every split. The QP family is an IPS- and DR-weighted BPR LightGCN baseline that is described in full
  in Appendix~A.
% source: results/v11/stageB.txt Sections 1 and 2
\end{description}

The comparison in C4 is not budget-matched. LEAF-LGN+DAX searches 3,840 configurations per split, whereas the QP family
alone searches 20. The fair comparator QP family+DAX adds the same DAX module to the QP family and searches 3,200
configurations. LEAF-LGN+DAX still exceeds it, by $+12.04$ [1.63] $\times10^{-3}$ nDCG@10, and we treat this
difference as the closer comparison.
% source: results/v11/stageB.txt Section 3 (configurations searched; fair baseline), Section 1 (n_configs)

We do not claim that DAX recovers an unbiased or causally identified item quality, since its factors are identified
only up to a scale; that LEAF-LGN+DAX outperforms SIM \\citep{pi2020sim} or TWIN \\citep{chang2023twin}, which were not run on the main setting; that the weight of
DAX is learned end to end; or that the results hold beyond one platform family and one training seed per split.
The deep lifelong baseline that we did run on the main setting is DIN \citep{zhou2018din}, which attends over the last $L$ events of each user's lifelong sequence. LEAF-LGN+DAX exceeds
DIN+DAX by $+5.73$ [1.64] $\times10^{-3}$ nDCG@10 on the same splits.
% source: results/v12/final12.txt Section 3b (LEAF-LGN+DAX - DIN+DAX); DIN definition in experiments/run_din.py

Section~\ref{sec:related} places the work among lifelong behaviour models, click models, exposure-aware recommenders,
occupancy and small-area estimators, and debiased graph recommenders. Section~\ref{sec:method} defines the LEAF
channels, DAX and their combination. Section~\ref{sec:protocol} gives the temporal protocol, the splits, the controls
and the tests. Section~\ref{sec:results} reports the headroom diagnosis, the main comparison, the backbone and ablation
analyses and the robustness checks. Section~\ref{sec:discussion} discusses why the multiplicative combination and the
lifelong channels help. Section~\ref{sec:limits} lists the limitations, and Section~\ref{sec:conclusion} concludes.
